\section{High-resolution model generation, study area, and observations}\label{sec:model}

\subsection{Automated high-resolution model generation}
The SWAT\textsuperscript{+} model is generated by the SWATGenX platform, which converts NHDPlus High Resolution hydrography, 1/3 arc-second 3DEP elevation, gSSURGO soils, and gridded land-use and climate forcings into an executable SWAT\textsuperscript{+} project under a documented preprocessing rule set. The decisive property for PFAS work is reach density: NHDPlus~HR resolves the stream network at roughly an order of magnitude finer reach spacing than the medium-resolution hydrography underlying earlier distributed PFAS models, so that ambient observations can be placed on hydrologically correct channels rather than on coarse aggregated subbasins. This resolution gain resides in network geometry and observation assignment, not in the resolution of meteorological forcing or PFAS source inventories, and is therefore a methodological advance specifically for the assignment of point observations and for the representation of sorbing constituents, whose load estimates are biased rather than merely noised by coarse discretization.

\subsection{Study area: the Rogue River and the Wolverine source}
We demonstrate the engine on the Rogue River watershed, a $606$~km\textsuperscript{2} tributary of the Grand River in Kent County, Michigan, gauged at Rockford (USGS~04118500) with a continuous daily streamflow record. The watershed is chosen deliberately. It contains the Wolverine World Wide tannery legacy and the House Street disposal site, where decades of leather-waterproofing waste and sludge disposal produced one of the most intensively investigated PFAS contaminated-site corridors in the United States, with a well-resolved downstream concentration signal in the lower Rogue. It is small enough to calibrate at high resolution---the generated model comprises $17{,}771$ HRUs, $590$ channels, $597$ routing units, and $6$ reservoirs---yet carries a gold-standard flow record and a dense ambient PFAS network. The clean point-source-to-river-gradient structure makes the Rogue a sharper test of source attribution than a watershed with diffuse, multi-source loading (Fig.~\ref{fig:study}).

\begin{figure}[ht]
\centering
\includegraphics[width=0.62\textwidth]{fig1_study_area.pdf}
\caption{The Rogue River watershed (USGS~04118500), Kent County, Michigan. The SWATGenX-generated SWAT\textsuperscript{+} network resolves $590$ channels from NHDPlus High Resolution. Black dots mark the $29$ EGLE ambient monitoring stations carrying quantified surface-water PFOS; the red star marks the Wolverine World Wide / House Street legacy source in the lower river. Inset: location within the Great Lakes basin.}
\label{fig:study}
\end{figure}

\subsection{Ambient PFOS observations and observation-to-channel assignment}
Ambient surface-water PFOS observations are drawn from the Michigan Department of Environment, Great Lakes, and Energy (EGLE) monitoring program: $31$ stations within the watershed, $29$ of which carry quantified (above-detection) surface-water PFOS and define the calibration set; the Rogue mainstem series among them resolves the downstream gradient toward the Wolverine corridor (rising from roughly $6$~ng~L\textsuperscript{-1} in the headwaters to $27$~ng~L\textsuperscript{-1} near the lower river). Because these observations are reported as per-station aggregate water concentrations rather than as daily time series, the calibration target is spatial---the pattern of representative concentration across stations---rather than temporal. Each station is assigned to its hydrologically correct model channel by a nearest-reach spatial join in an equal-area projection, mirroring the assignment SWATGenX uses for streamflow gauges but using proximity rather than drainage-area matching because concentration stations, unlike flow gauges, are not constrained to integrate a known upstream area. The $31$ stations map to $22$ distinct channels at a median snap distance of a few meters; the $29$ with quantified PFOS fall on $20$ of those reaches, which constitute the calibration target (several stations share a reach, so metrics are reported per reach using the station-maximum observed concentration---the source-relevant value). This high-resolution assignment is what allows the spatial gradient---and therefore the source signal---to be compared reach by reach, and it is not available at the coarser discretization of prior distributed PFAS models of comparable watersheds.

\subsection{Soil-PFAS initialization}
Initial soil PFAS pools and the Freundlich and Langmuir coefficients are assigned per HRU from soil class and land use: gSSURGO-derived soil series are mapped to the sorption-parameter classes of the prior distributed model \citep{rafiei2023}, and land use sets a background initial soil concentration, converted to a per-layer mass using bulk density and layer thickness. This national, automatable initialization replaces the measured per-site soil inventory that is available only for data-rich legacy basins, and the overall magnitude of the resulting background is the single quantity adjusted in calibration (Section~4). The resulting per-HRU fields (Fig.~\ref{fig:soil}) show the spatial heterogeneity of the initial soil pool and the discrete, soil-class-controlled sorption parameters that the engine reads.

\begin{figure}[ht]
\centering
\includegraphics[width=\textwidth]{fig4_soil_params.pdf}
\caption{Per-HRU soil-PFAS initialization across the $17{,}771$ Rogue HRUs: (a) initial surface-layer soil PFOS pool (continuous, NaturalBreaks classes); (b) Freundlich sorption coefficient $k_f$ and (c) Freundlich exponent $n$, both discrete soil-class fields. Colour map: batlow (perceptually uniform, colour-blind safe). The three panels share the same spatial extent; the scale bar and north arrow on panel~(a) apply to all.}
\label{fig:soil}
\end{figure}
